\section*{致谢}
我们感谢 Jean Zinn--Justin 就 Langevin 方程的重正化与我们进行的
有益讨论，以及他鼓励我们直接处理流方程，
而不是把后者看作 Langevin 方程的某种极限。

\appendix

\section{记号约定}\label{app:A}

$\SUn$ 的李代数 $\sun$
可以等同于所有反厄米、无迹的 $N\times N$ 矩阵构成的线性空间。
取这样一组矩阵的基 $T^a$，$a=1,\ldots,N^2-1$，
元素 $X\in\sun$ 便可写作 $X=X^aT^a$，其中分量 $X^a$ 为实数
（重复的群指标自动求和）。对易关系
\begin{equation}
  [T^a,T^b]=f^{abc}T^c
  \label{eq:A.1}
\end{equation}
中的结构常数 $f^{abc}$ 在归一化条件
\begin{equation}
  \tr\{T^aT^b\}=-\frac{1}{2}\delta^{ab}
  \label{eq:A.2}
\end{equation}
被施加时为实数且对指标全反对称。
此外，$f^{acd}f^{bcd}=N\delta^{ab}$。

连续理论中的规范场取值于规范群的李代数。
洛伦兹指标 $\mu,\nu,\ldots$ 在成对出现时自动求和。
时空度规取为欧几里得的。特别地，对任意动量 $p$ 有 $p^2=p_{\mu}p_{\mu}$，
且 $\delta_{\mu\mu}=D$。

\section{流顶点}\label{app:B}

流顶点 $\vbulk{2,0}$、$\vbulk{3,0}$ 与 $\vbulk{1,1}$
分别通过式 \eqref{eq:2.13} 与 \eqref{eq:6.12} 定义。
它们的显式表达式为
\begin{align}
  \vbulk{2,0}(p,q,r)^{abc}_{\mu\nu\rho}=
  if^{abc}\bigl\{(r-q)_{\mu}\delta_{\nu\rho}
  +2q_{\rho}\delta_{\mu\nu}-2r_{\nu}\delta_{\mu\rho}
  \notag\\[2.5ex]
  +(\alpha_0-1)
  (q_{\nu}\delta_{\mu\rho}-r_{\rho}\delta_{\mu\nu})\bigr\},
  \label{eq:B.1}\\[3.5ex]
  \vbulk{3,0}(p,q,r,s)^{abcd}_{\mu\nu\rho\sigma}=
  f^{abe}f^{cde}
  (\delta_{\mu\sigma}\delta_{\nu\rho}-\delta_{\mu\rho}\delta_{\sigma\nu})
  \notag\\[2.5ex]
  \quad
  +f^{ade}f^{bce}
  (\delta_{\mu\rho}\delta_{\sigma\nu}-\delta_{\mu\nu}\delta_{\rho\sigma})+
  f^{ace}f^{dbe}
  (\delta_{\mu\nu}\delta_{\rho\sigma}-\delta_{\mu\sigma}\delta_{\nu\rho}),
  \label{eq:B.2}\\[3.5ex]
  \vbulk{1,1}(p,q,r)^{abc}_{\mu}=\alpha_0if^{abc}r_{\mu}.
  \label{eq:B.3}
\end{align}
在费曼图中，这些公式里的动量取为入射线的动量。

\section{图~\ref{fig:brsdiag} 中各图的计算}\label{app:C}

图~\ref{fig:brsdiag} 中的各图对 BRS 恒等式 \eqref{eq:6.23}
中关联函数的傅里叶变换的贡献形如
\begin{equation}
   -{g_0^4if^{abc}\over \lambda_0 p^2q^2r^2}\mathcal{C}_{\mu,n},
   \label{eq:C.1}
\end{equation}
其中 $n$ 是图的编号。具体地，对出现在方程左端的各图有
\begin{align}
   \mathcal{C}_{\mu,1}=
   \frac{1}{2}p^2(q-r)_{\mu}\rme^{-t\alpha_0(q^2+r^2)}+
   \frac{1}{2}(q^2-r^2)p_{\mu}\rme^{-t\alpha_0p^2}
   -\mathcal{C}_{\mu,2},
   \label{eq:C.2}\\[3pt]
   \mathcal{C}_{\mu,2}=
   \frac{1}{2}\bigl\{p^2(q-r)_{\mu}+(q^2-r^2)p_{\mu}\bigr\}\rme^{-tp^2},
   \label{eq:C.3}
\end{align}
而对右端的各图有
\begin{align}
   \mathcal{C}_{\mu,3}=-p^2r_{\mu}\rme^{-t\alpha_0(q^2+r^2)},
   \label{eq:C.4}\\[3pt]
   \mathcal{C}_{\mu,4}=-\frac{1}{2}p^2p_{\mu}\bigl\{
   \rme^{-t\alpha_0(q^2+r^2)}-\rme^{-t\alpha_0p^2}\bigr\},
   \label{eq:C.5}\\[3pt]
   \mathcal{C}_{\mu,5}=-(qr+r^2)p_{\mu}\rme^{-t\alpha_0p^2}.
   \label{eq:C.6}
\end{align}
BRS 对称性要求
\begin{equation}
   \mathcal{C}_{\mu,1}+\mathcal{C}_{\mu,2}=
   \mathcal{C}_{\mu,3}+\mathcal{C}_{\mu,4}+\mathcal{C}_{\mu,5},
   \label{eq:C.7}
\end{equation}
事实确实如此。

\begin{thebibliography}{7}

\bibitem{WilsonFlow}
M. L\"uscher,
{\itshape Properties and uses of the Wilson flow in lattice QCD},
JHEP 1008 (2010) 071

\bibitem{Villasimius}
M. L\"uscher,
{\itshape Topology, the Wilson flow and the HMC algorithm},
XXVIII International Symposium on Lattice Field Theory, June 14-19 2010, Villasimius, Italy,
PoS(Lattice 2010)015

\bibitem{Zinn}
J. Zinn--Justin,
{\itshape Renormalization and stochastic quantization},
Nucl. Phys. B275 [FS17] (1986) 135

\bibitem{ZinnZwanziger}
J. Zinn--Justin, D. Zwanziger,
{\itshape Ward identities for the stochastic quantization of gauge fields},
Nucl. Phys. B295 [FS21] (1988) 297

\bibitem{BRSI}
C. Becchi, A. Rouet and R. Stora,
{\itshape Renormalization of the abelian Higgs-Kibble model},
Comm. Math. Phys. 42 (1975) 127

\bibitem{BRSII}
C. Becchi, A. Rouet and R. Stora,
{\itshape Renormalization of gauge theories},
Ann. Phys. (NY) 98 (1976) 287

\bibitem{Symanzik}
K. Symanzik,
{\itshape Schr\"odinger representation and Casimir effect in
renormalizable quantum field theory},
Nucl. Phys. B190 [FS3] (1981) 1

\end{thebibliography}
